\section{Trabalhos Relacionados}

O fenômeno investigado neste trabalho -- aumento de salinidade em efluente de estação de tratamento associado a medidas de conservação de água -- foi estudado por Schwabe et al. \cite{schwabe2020unintended} em 34 estações de tratamento do sul da Califórnia, analisando vazão e salinidade mensais de 2013 a 2017 (período de uma das secas mais intensas já registradas na região). Os autores relatam que a conservação de água durante a seca reduziu o volume de esgoto que chega às estações e aumentou sua salinidade -- o mesmo mecanismo (mesma massa de sais, menor volume de água) usado neste trabalho para interpretar a tendência de alta de TDS encontrada na LAGWRP (Seção de Conclusão). Não foi possível acessar o texto completo do artigo nesta sessão (acesso pago); a citação aqui se apoia no resumo divulgado pela instituição dos autores (University of California, Riverside), que confirma o desenho do estudo e a direção do efeito, mas não traz a magnitude exata do aumento de salinidade -- por isso este trabalho não reproduz nenhum número específico desse estudo, só a direção qualitativa do efeito.

Wolfand et al. \cite{wolfand2022dilution} modelam, via EPA SWMM, o impacto de cenários de reúso de água (0\%, 50\% e 100\%) e de desvio de vazão de galerias pluviais sobre a qualidade da água na bacia do rio Los Angeles -- a mesma bacia onde a LAGWRP descarrega seu efluente (pontos de monitoramento R-4, R-7, RSW-650 e RSW-654 do dataset eSMR usado neste trabalho, tratados aqui apenas como contexto). Os autores relatam que a concentração de sólidos dissolvidos totais (TDS), sólidos suspensos totais e alguns metais aumenta com o reúso de água, mesmo quando a carga total (massa) desses poluentes diminui -- reforçando, numa bacia hidrográfica real e não apenas em nível de estação individual, o mesmo mecanismo de concentração por redução de volume investigado neste trabalho. Como no caso anterior, o acesso ao texto completo foi bloqueado nesta sessão; a citação se apoia em resumos de terceiros (ASCE, sociedade profissional de engenharia civil) que confirmam o desenho do estudo e a direção do efeito para TDS, sem fornecer a magnitude exata do aumento -- nenhum valor numérico deste estudo é reproduzido aqui por esse motivo.

Porse et al. \cite{porse2023adapting}, diferentemente dos dois trabalhos anteriores, teve seu texto completo obtido e lido na íntegra nesta sessão (artigo de acesso aberto, CC BY-NC-ND). Os autores combinam revisão de literatura histórica com pesquisa e entrevistas junto a gestores de sistemas de coleta, tratamento e reúso de esgoto em toda a Califórnia (81 respostas de estações de tratamento, cobrindo 65 agências e 84 instalações distintas), financiados pelo California State Water Resources Control Board. O achado mais diretamente relevante para este trabalho: para 45 estações de tratamento do sul urbano da Califórnia, as exigências de conservação de água ligadas à seca de 2014-2016 resultaram em queda de 6-10\% na vazão de efluente (controlando por outros fatores), e essas mesmas ações de conservação aumentaram o TDS do efluente em 5-12\% -- uma cifra estadual, percentual e verificada (diferente do valor absoluto "+50 mg/L" atribuído a Wolfand et al., que segue não confirmado por falta de acesso ao texto completo daquele estudo). Os autores relatam ainda que 63\% das estações de tratamento pesquisadas relataram desafios operacionais ligados a baixa vazão desde 2011 -- o mesmo ano de início da série temporal da LAGWRP usada neste trabalho -- o que reforça que o fenômeno de queda de vazão pós-2011 é um padrão reconhecido estadualmente pela própria indústria de saneamento californiana, e não uma particularidade isolada da estação aqui analisada.

Do ponto de vista metodológico, o tratamento dos valores não detectados do BOD (Seção de Metodologia) segue a avaliação sistemática de Antweiler e Taylor \cite{antweiler2008evaluation}, que comparam substituição por MDL/2, máxima verossimilhança, Kaplan-Meier e ROS (\textit{Regression on Order Statistics}) usando dados ambientais reais com valor verdadeiro conhecido.

O estudo de maior influência sobre o enquadramento final deste trabalho, no entanto, é o relatório técnico da Southern California Salinity Coalition \cite{scsc2018tds}, preparado pela consultoria Daniel B. Stephens \& Associates para 26 estações de tratamento do sul da Califórnia -- o único dos quatro trabalhos citados nesta seção cujo texto completo foi efetivamente lido nesta sessão (via conversão para texto, \texttt{material\_apoio\_referencias.md} Tema 9), não apenas resumo de terceiros. Usando decomposição de importância relativa (método LMG, pacote R \texttt{relaimpo}), os autores atribuem a maior parte da variabilidade do TDS de influente (~88\%, contra ~12\% do consumo per capita) ao TDS da própria água de abastecimento, que varia com a origem da água importada (\textit{State Water Project}, \textit{Colorado River Aqueduct}) em resposta a ciclos de seca -- não à conservação local em si. Esse achado reenquadra o mecanismo de Schwabe et al.\ \cite{schwabe2020unintended} como um efeito real, porém secundário, e foi a motivação direta para testar o PDSI (Índice de Severidade de Seca de Palmer) como variável explicativa neste trabalho (Seção de Resultados, ``PDSI como teste independente do reenquadramento cíclico''), e para adotar o método LMG na decomposição PDSI-vazão do WRTDS (Seção de Metodologia).
